% =================================================================
%  30-Month Doctoral Report
%  Tengjiao Sun -- Vision, Learning and Control, ECS, Southampton
%
%  This file is self-contained and does NOT touch main.tex (9-month)
%  or 18_months.tex (18-month).  Compile with:
%      pdflatex 30_months && bibtex 30_months && pdflatex 30_months x2
% =================================================================
\documentclass[11pt,a4paper]{report}

\usepackage[utf8]{inputenc}
\usepackage[T1]{fontenc}
\usepackage[margin=1in]{geometry}
\usepackage{amsmath,amssymb}
\usepackage{graphicx}
\usepackage{booktabs}
\usepackage{tabularx}
\usepackage{enumitem}
\usepackage{fancyhdr}
\usepackage{titlesec}
\usepackage{xcolor}
\usepackage{pgfgantt}
\usepackage[hidelinks]{hyperref}

% ---- running header in the style of the VLC report template ------
\pagestyle{fancy}
\fancyhf{}
\renewcommand{\headrulewidth}{0pt}
\fancyhead[L]{\small\itshape\nouppercase{\leftmark}}
\fancyhead[R]{\small\thepage}
\fancyfoot{}
\renewcommand{\chaptermark}[1]{\markboth{Chapter \thechapter.\ #1}{}}

% ---- compact chapter headings ------------------------------------
\titleformat{\chapter}[display]
  {\normalfont\huge\bfseries}{\chaptertitlename\ \thechapter}{10pt}{\Huge}
\titlespacing*{\chapter}{0pt}{-20pt}{25pt}

\setlist[itemize]{topsep=2pt,itemsep=1pt,parsep=0pt,leftmargin=*}
\setlist[enumerate]{topsep=2pt,itemsep=1pt,parsep=0pt,leftmargin=*}

% ---- marker for content only the author can supply ---------------
% Set \suppressnotes to true before submitting to hide every marker.
\newif\ifsuppressnotes\suppressnotesfalse
\newcommand{\needsinput}[1]{%
  \ifsuppressnotes\else\textcolor{orange!85!black}{\textbf{[fill in: }#1\textbf{]}}\fi}

\newcommand{\workitem}[4]{%
  \noindent\textbf{Relation to thesis:} #1\par\smallskip
  \noindent\textbf{Goal:} #2\par\smallskip
  \noindent\textbf{Hypothesis:} #3\par\smallskip
  \noindent\textbf{Tasks:} #4\par}

\begin{document}

% =================================================================
\begin{titlepage}
\centering
\vspace*{1.5cm}
{\large\scshape University of Southampton}\\[1.2cm]
{\scshape Faculty of Engineering and Physical Sciences}\\
{\scshape Electronics and Computer Science}\\[2.5cm]
{\LARGE\bfseries Controllable and Generalisable\\[0.25cm]
3D Human Motion Generation\\[0.25cm]
Using Multimodal Approaches\par}
\vspace{2.5cm}
\begin{minipage}{0.45\textwidth}
\begin{flushleft}\large
\emph{Author:}\\
Tengjiao Sun
\end{flushleft}
\end{minipage}
\begin{minipage}{0.45\textwidth}
\begin{flushright}\large
\emph{Supervisors:}\\
Dr Hansung Kim\\
Dr Zhiwu Huang
\end{flushright}
\end{minipage}\\[3cm]
{\itshape A report submitted in fulfillment of the requirements for the
30 month report\\ in the Vision, Learning and Control research group}\\[2cm]
{\large September, 2026}
\vfill
\end{titlepage}

% =================================================================
\chapter{Thesis Structure}

The following table details the main focus of each chapter of the final thesis.
Chapters~3--6 correspond to work that is complete and either published or under
review; Chapter~7 corresponds to work in progress and is described in
Chapter~3 of this report. A more detailed structure, including subsections,
follows the table.

\begin{table}[h]
\centering\small
\begin{tabularx}{\textwidth}{@{}l X X@{}}
\toprule
\textbf{Chapter} & \textbf{Research Focus} & \textbf{Key Contribution}\\
\midrule
Chapter 1 & Introduction & --\\
Chapter 2 & Literature Review & --\\
Chapter 3 & Shared Conditional Representation for Text and Motion (MCRE)
          & A joint text--motion condition space admitting linear semantic
            editing directions\\
Chapter 4 & Hierarchical One-Pass Autoregressive Generation (Hi-RQCT, MOGO)
          & Real-time, infinite-length generation via scale-adaptive residual
            quantization and one-pass causal decoding\\
Chapter 5 & Part-Tokenized Motion and Attention-Side Control (PartVQ,
            KV-Control)
          & An anatomy-aligned token substrate, and sub-centimetre trajectory
            control injected into a frozen prior as attention memory\\
Chapter 6 & Topology-Agnostic Motion Generation (TopX)
          & Generation over arbitrary skeletons without a fixed reference
            topology\\
Chapter 7 & Multi-Subtask Motion Generation & [TBD]\\
Chapter 8 & Conclusion \& Future Work & --\\
\bottomrule
\end{tabularx}
\caption{Summary of thesis structure and research contributions}
\label{tab:structure}
\end{table}

\section{Section structure}

\begin{enumerate}[leftmargin=1.4em]
\item \textbf{Introduction}
  \begin{enumerate}[leftmargin=1.4em]
  \item Background
  \item Problem statement
  \item Scope
  \item Contributions
  \end{enumerate}

\item \textbf{Literature Review: Controllability and Generality in Human
  Motion Generation}
  \begin{enumerate}[leftmargin=1.4em]
  \item Text-to-motion: VAEs, tokenization and diffusion
  \item Motion representation and quantization
  \item Spatial control, editing and composition
  \item Skeleton topology, retargeting and transfer
  \item Synopsis of literature
  \end{enumerate}

\item \textbf{Shared Conditional Representation for Text and Motion}\\
  This chapter builds the conditioning interface used throughout the thesis.
  A lightweight motion adapter is aligned to a frozen pre-trained text encoder
  so that captions and motions occupy one space, in which semantic edits are
  displacement vectors and multiple conditions compose by addition. The chapter
  also establishes the granularity limit that the rest of the thesis removes:
  a clip-level embedding says what a motion is, not where or when any part of
  it happens.

\item \textbf{Hierarchical One-Pass Autoregressive Generation}
  \begin{enumerate}[leftmargin=1.4em]
  \item Overview
  \item Scale-adaptive residual quantization and cross-level decorrelation
  \item One-pass hierarchical causal decoding
  \item Infinite-length generation and mid-stream prompt switching
  \item Results and evaluation
  \item Conclusion
  \end{enumerate}

\item \textbf{Part-Tokenized Motion and Attention-Side Control}
  \begin{enumerate}[leftmargin=1.4em]
  \item Overview
  \item PartVQ: anatomy-aligned part codebooks
  \item T-Concat: frame--part tokens as attention-addressable sites
  \item KV-Control: control-conditioned key/value injection
  \item Results and evaluation
  \item Conclusion
  \end{enumerate}

\item \textbf{Topology-Agnostic Motion Generation}
  \begin{enumerate}[leftmargin=1.4em]
  \item Overview
  \item Methodology
  \item Results and evaluation
  \item Out-of-distribution specialisation
  \item Conclusion
  \end{enumerate}

\item \textbf{Multi-Subtask Motion Generation}
  \begin{enumerate}[leftmargin=1.4em]
  \item Overview
  \item Methodology
  \item Results and evaluation
  \item Conclusion
  \end{enumerate}

\item \textbf{Conclusion}
  \begin{enumerate}[leftmargin=1.4em]
  \item Summary
  \item Impact and applications
  \item Limitations
  \item Future work
  \end{enumerate}
\end{enumerate}

% =================================================================
\chapter{Summary of Work}

\section{Research narrative}

This thesis builds a text-to-motion system that is controllable at the level
practitioners actually work at --- a named joint, a sketched path, a specific
interval --- and general enough to survive a change of body or a change of
task. The work reaches this through a single technical commitment: motion is
represented as a \emph{part-tokenized substrate}, and control is expressed as
an operation on that substrate rather than as a stronger condition on a
generator.

The first work, MCRE~\cite{sun2024mcre}, built the conditioning interface. A
lightweight motion adapter trained against a frozen text
encoder~\cite{radford2021learning} places captions and motions in one space, so
that semantic edits become displacement vectors and conditions compose by
addition. This is the interface every later chapter inherits. It also fixes the
granularity of what it can express: one vector per clip, which says what a
motion is but not where or when any part of it happens.

The second work made generation fast and unbounded. Diffusion text-to-motion
models~\cite{tevet2022human,chen2023executing,zhang2024motiondiffuse} sample
iteratively and train at a fixed horizon, which rules out interactive use.
Hi-RQCT~\cite{sun2025hirqct} and then MOGO~\cite{sun2025mogo} discretise motion
with a scale-adaptive residual vector
quantizer~\cite{vandenoord2017neural,zeghidour2022soundstream} and decode every
quantization level in a single causal forward pass, giving real-time synthesis,
unbounded sequence length and mid-stream prompt switching. MOGO was published at
AAAI 2026. This established the discrete substrate the thesis works on, and
sharpened the control problem: a token grid is fast, but a flat grid of
whole-body codes offers nothing to address when a user wants one hand in one
place.

The third work is the centre of the thesis. KV-Control~\cite{sun2026kvcontrol}
makes the substrate addressable and then uses that addressability to control a
model it never retrains. \textbf{PartVQ} learns anatomy-aligned part codebooks,
so each token names a body part at a frame instead of a whole pose.
\textbf{T-Concat} exposes each frame--part token as an attention-addressable
site. \textbf{KV-Control} then injects control-conditioned key/value memories at
every self-attention layer of a frozen masked text-to-motion transformer,
leaving the query stream, the text cross-attention, the feed-forward blocks and
all backbone weights untouched. Geometric constraints become memory that
attention can retrieve, rather than a global pose token bolted onto the input or
a penalty applied at the output. The adapter trains only injection parameters on
top of a shared trajectory encoder, and tracks root and multi-joint trajectories
to sub-centimetre accuracy while retaining the text-conditioned motion quality
of the prior. This is the point at which the thesis delivers precise spatial
control without the usual price: no duplicated generator branch, no test-time
optimisation, no degradation of the pretrained prior.

PartVQ has since become the substrate other work in the group builds on. In
MoGeFlow~\cite{fang2026mogeflow} the part codebooks are shown to carry
measurable, decoder-causal kinematic geometry, which licenses generation as a
continuous flow through codebook space rather than categorical index
prediction; in AnchorRoute~\cite{fang2026anchorroute} the same addressable
tokens carry sparse anchors both as generation-time conditioning and as
residuals that route interval-local refinement. That two further systems are
built on it is the strongest available evidence that the substrate, not any one
controller, is the durable contribution.

Two limits remain, and both are about generality rather than precision. First,
all of the above assumes one body: HumanML3D~\cite{guo2022generating},
KIT-ML~\cite{plappert2016kitkit} and AMASS~\cite{mahmood2019amass} share a
single SMPL~\cite{loper2015smpl} skeleton, and PartVQ's part decomposition is
defined on it, so transfer to another character runs through
retargeting~\cite{aberman2020skeletonaware} and its correspondence assumption.
TopX removes this: the skeleton becomes an input to the model rather than a
property of the dataset. Second, all of the above executes one instruction.
Real requests are ordered and composite, and concatenating independently
generated clips leaves seams and drifting
style~\cite{shafir2024priormdm,barquero2024flowmdm}. The multi-subtask work
addresses this directly, and the part-tokenized substrate is what makes it
tractable: subtask boundaries can be written into a sequence whose tokens are
already addressable per part and per frame.

These two threads define the remaining period and are set out in Chapter~3.

\section{Completed contributions}

\subsection{Shared conditional representation for text and motion}

Text-to-motion models must reconcile modalities that differ in dimensionality,
length and structure. The standard approach conditions a generator on an
embedding from an encoder that has never seen motion, which couples the two
representations loosely and makes anything other than one-shot generation ---
editing, interpolation, composition --- awkward to express.

This work trained a lightweight motion adapter against a frozen pre-trained
text encoder, mapping a motion and its caption to nearby points in one
conditional space. Because the space is shared and the objective is
contrastive, semantic relationships between captions transfer to motions:
changing the speed, direction or affect of a movement is a displacement in the
conditional space rather than a new prompt, and several conditions combine by
composition in the same space.

The result is the conditioning interface reused by every subsequent chapter,
and a precise statement of what clip-level conditioning cannot do. That
statement sets the agenda for Chapters~4 and 5.

Published at the ECCV 2024 Workshop on Foundation Models for 3D
Humans~\cite{sun2024mcre}.

\subsection{Hierarchical one-pass autoregressive generation}

Tokenization-based motion models~\cite{chuan2022tm2t,guo2024momask} decode
faster than diffusion, but residual quantization schemes generally require one
decoding pass per quantization level, and error accumulates over long
autoregressive rollouts.

This work contributed two components. MoSA-VQ is a residual vector quantizer
with learnable per-level scaling and a cross-level decorrelation objective; the
decorrelation term minimises the batch-wise covariance between a quantized
vector and the residual passed to the next level, enforcing approximate
statistical orthogonality across the hierarchy so each level encodes what the
previous levels did not. This improves codebook utilisation and yields more
compact codes under a fixed budget. The RQHC-Transformer is a causal decoder
whose structure mirrors the quantization hierarchy and which emits all levels in
one forward pass, removing per-level decoding cost while preserving temporal
coherence.

Together these give near real-time decoding, generation of unbounded length,
and injection of a new text prompt at an arbitrary point in the sequence --- the
property that makes the system usable for avatar control and teleoperation
rather than offline clip synthesis. A textual condition alignment module
normalises free-form input into structured instructions before generation,
improving robustness on prompts unlike those seen in training.

Ablations confirmed both components. Removing the decorrelation objective
degraded generation FID and retrieval accuracy consistently across metrics.
Varying codebook capacity showed that the configuration minimising
reconstruction error is not the one maximising generation quality, which
established that quantization error here is structured rather than uniform
noise --- the observation that motivated moving to an anatomy-aligned
decomposition in the next work.

The decoder was published at CVMP 2025 as Hi-RQCT~\cite{sun2025hirqct}; the full
system was published at AAAI 2026 as MOGO~\cite{sun2025mogo}.

\subsection{Part-tokenized motion and attention-side control}

Text-conditioned models synthesise plausible motion from prompts, but animation
and embodied-agent workflows do not stop at text: a character must follow a
sketched root path, reach an end-effector target, or satisfy a multi-joint
trajectory while keeping the gait and intent the language described. Existing
trajectory controllers pay for this either by duplicating large parts of the
generator to regain per-layer access, or by pushing the cost to test-time
optimisation.

This work provides a compact attention-side control interface for frozen masked
text-to-motion transformers, built on three components:

\begin{itemize}
\item \textbf{PartVQ} learns anatomy-aligned part codebooks, so a token denotes
  a body part at a frame rather than a whole pose. This is what makes the
  token grid addressable at the granularity control actually needs.
\item \textbf{T-Concat} exposes each frame--part token as a site that attention
  can address directly.
\item \textbf{KV-Control} injects control-conditioned key/value memories at
  every self-attention layer, while preserving the pretrained query stream, the
  text cross-attention, the feed-forward blocks and all backbone weights. Only
  injection parameters on top of a shared trajectory encoder are trained.
\end{itemize}

The reframing is that geometric constraints are made available \emph{as memory
inside self-attention} rather than delivered through a global pose token or
enforced at the output. Trajectory conditioning becomes lightweight memory
retrieval. The resulting adapter tracks root and multi-joint constraints to
sub-centimetre accuracy under the inherited refinement protocol while retaining
the text-conditioned motion quality of the frozen prior, and it is small,
precise and inspectable --- the control path is a named set of injected
memories rather than a re-trained generator.

The paper is under review at AAAI.
\needsinput{confirm the current title under which this was submitted, since the
report refers to it as KV-Control throughout}

\subsection{Topology-agnostic motion generation}

Every model above is trained on one skeleton. The standard benchmarks derive
from SMPL~\cite{loper2015smpl}, and PartVQ's anatomical decomposition is defined
against it, so animating a structurally different character runs through
retargeting~\cite{aberman2020skeletonaware}, which assumes a joint
correspondence between source and target. For morphologies with different limb
counts, proportions or non-humanoid structure, that correspondence does not
exist and the assumption fails silently.

TopX generates motion for skeletons of arbitrary topology with no fixed
reference structure: the skeleton is an input to the model rather than a
property of the dataset.
\needsinput{the methodology --- how the skeleton is encoded and consumed, what
the training corpus is, and what separates TopX from
AnyTop~\cite{gilles2024anytop} and from retargeting-based approaches}

\needsinput{the headline quantitative result, ideally performance on a skeleton
class held out from training}

The paper is complete and in submission to ICLR 2027. Behaviour on skeletons
outside the training distribution is the subject of Section~3.1.

\subsection{Contributing-author work}

Three further motion systems were developed with colleagues in the group, each
building on components contributed here.

\begin{itemize}
\item \textbf{MotionDuet}~\cite{zhang2026motionduet}, published at CVPR 2026,
  aligns motion generation with the distribution of video-derived
  representations. Video cues from a pretrained model ground low-level dynamics
  while text supplies semantic intent; DUET fuses video-informed cues into the
  motion latent space, a DASH loss aligns generated trajectories with the
  distributional and structural statistics of video features, and an
  auto-guidance mechanism balances the two conditions using a weakened copy of
  the model. My contribution was to the dual-conditioning design and the
  training and evaluation of the fused model.

\item \textbf{MoGeFlow}~\cite{fang2026mogeflow} shows that PartVQ's
  group-specific codes carry measurable, non-random and decoder-causal
  kinematic geometry, and generates by learning a text-conditioned continuous
  flow over code embeddings rather than predicting categorical indices. It sets
  state of the art in R-Precision on HumanML3D and KIT-ML. My contribution was
  the PartVQ substrate it is built on and the codebook-geometry diagnostics.

\item \textbf{AnchorRoute}~\cite{fang2026anchorroute} uses sparse anchors as a
  single scaffold for both generation and refinement: anchors are injected into
  a frozen masked diffusion prior as condition memory, and the same anchors are
  then read as residuals whose timestamps define refinement intervals. My
  contribution was the attention-side injection mechanism carried over from
  KV-Control.

\item \textbf{UMO}~\cite{zhan2026umo} addresses unified sparse motion modelling
  for real-time co-speech avatars, to which I contributed on the sparse motion
  representation.
\end{itemize}

\section{Publications}

\begin{itemize}
\item \textbf{ECCV 2024 Workshop} --- \emph{MCRE: Multimodal Conditional
  Representation and Editing for Text-Motion Generation}. First author.
  Established the joint text--motion conditional space used throughout the
  thesis.

\item \textbf{CVMP 2025} --- \emph{Hi-RQCT: Hierarchical Residual-Quantized
  Causal Transformer for High-Quality 3D Human Motion Generation}. Joint first
  author. The hierarchical causal decoder later extended into MOGO.

\item \textbf{AAAI 2026} --- \emph{MOGO: Residual Quantized Hierarchical Causal
  Transformer for Real-Time and Infinite-Length 3D Human Motion Generation}.
  Joint first author. Real-time decoding, unbounded length and mid-stream
  prompt switching.

\item \textbf{CVPR 2026} --- \emph{MotionDuet: Dual-Conditioned 3D Human Motion
  Generation with Video-Regularized Text Learning}. Second author.

\item \textbf{AAAI, under review} --- \emph{KV-Control: Parameter-Efficient K/V
  Injection for Trajectory-Controlled Text-to-Motion}. First author. PartVQ,
  T-Concat and attention-side control injection.

\item \textbf{ICLR 2027, in submission} --- \emph{TopX}. First author.
  Generation over arbitrary skeleton topologies.

\item \textbf{Contributing author, 2026} ---
  \emph{MoGeFlow}~\cite{fang2026mogeflow},
  \emph{AnchorRoute}~\cite{fang2026anchorroute} and
  \emph{UMO}~\cite{zhan2026umo}.
\end{itemize}

% =================================================================
\chapter{Plan of Unfinished Work}

Two threads remain, addressing the two generality limits identified at the end
of Section~2.1: dependence on a single body, and dependence on a single
instruction.

\section{Out-of-distribution specialisation for topology-agnostic generation}

\workitem{%
TopX makes the skeleton an input rather than a dataset property. This work
establishes what that buys on skeletons far from the training distribution, and
how cheaply the model can be specialised to a new one. It completes Chapter~6.}{%
To characterise and then reduce the degradation of topology-agnostic generation
under out-of-distribution skeletons --- morphologies, limb counts and
proportions absent from training --- and to determine how much target-domain
data is needed to specialise the model to a new topology.}{%
A topology-agnostic model learns two separable things: motion dynamics, which
transfer across morphology, and a structural prior over skeletons, which does
not. Out-of-distribution failure is therefore concentrated in the structural
component, is visible in internal representations before it is visible in output
quality, and can be repaired with a small amount of target data by adapting only
that component --- with generative replay from the base model preventing the
adaptation from erasing the dynamics.}{%
\begin{enumerate}
\item Build an out-of-distribution evaluation protocol: held-out skeleton
  classes at controlled distance from the training distribution, with quality
  reported as a function of that distance rather than as one aggregate
  number~\cite{hendrycks2021ood}.
\item Locate the degradation by separating dynamics and structural components
  in the representation, and test whether failure is predictable from internal
  representations ahead of output quality.
\item Build the specialisation procedure: parameter-efficient adaptation of the
  structural component with generative replay, and measure the data--quality
  trade-off curve for a new topology.
\item Evaluate against retargeting baselines~\cite{aberman2020skeletonaware} and
  AnyTop~\cite{gilles2024anytop} on the held-out classes.
  \needsinput{which OOD skeleton sets are available, and whether a non-humanoid
  corpus is in hand}
\end{enumerate}}

\section{Multi-subtask motion generation}

\workitem{%
Every model in the thesis executes a single instruction. This work extends the
conditioning interface of Chapter~3 to ordered, composite instructions, and is
the last piece required for the thesis claim of control that is both precise and
general.}{%
To generate long-horizon motion from an ordered set of subtasks such that
transitions are generated rather than stitched, style and identity persist
across subtask boundaries, and the per-part control of Chapter~5 remains
available within any individual subtask.}{%
Existing approaches treat composition as a post-hoc problem: clips are generated
independently and blended~\cite{shafir2024priormdm,barquero2024flowmdm}, so the
transition is a smoothing artefact rather than something the model has learned.
Writing subtask boundaries into the sequence the model decodes makes the
transition part of the modelled distribution, improving both boundary quality
and global style consistency over blending. The part-tokenized substrate of
Chapter~5 is what makes this practical: tokens are already addressable per part
and per frame, so a boundary can be expressed without leaving the
representation, and MOGO's causal decoder already accepts a condition change
mid-sequence.}{%
\begin{enumerate}
\item Fix the subtask representation: how a boundary is encoded, and whether
  ordering is supplied as a hard schedule or inferred.
  \needsinput{the current design, and which of the two the model assumes}
\item Assemble training data with subtask annotations, from compositional
  corpora~\cite{lee2023multiact,zhang2023finemogen} or by segmenting long
  sequences. \needsinput{the dataset in use}
\item Train the compositional decoder on the part-tokenized backbone and verify
  that per-part controllability survives composition.
\item Evaluate transition quality, long-horizon drift and style consistency
  against blending-based baselines, and confirm that KV-Control injection still
  applies within a subtask.
\end{enumerate}}

\section{Timeline}

\begin{figure}[h]
\centering
\begin{ganttchart}[
  x unit=0.63cm,
  y unit title=0.55cm,
  y unit chart=0.55cm,
  vgrid, hgrid,
  title label font=\tiny,
  bar label font=\scriptsize,
  group label font=\scriptsize\bfseries,
  milestone label font=\scriptsize,
  bar/.append style={fill=blue!40},
  group/.append style={fill=black!65},
  milestone/.append style={fill=orange!80},
  bar height=.5,
  group right shift=0,
  group left shift=0
]{1}{13}
\gantttitle{2026}{4}
\gantttitle{2027}{9}\\
\gantttitlelist{9,...,12}{1}
\gantttitlelist{1,...,9}{1}\\

\ganttgroup{TopX / OOD}{1}{7}\\
\ganttbar{OOD evaluation protocol}{1}{3}\\
\ganttbar{Degradation diagnosis}{3}{5}\\
\ganttbar{Specialisation + replay}{4}{7}\\
\ganttmilestone{ICLR decision / resubmission}{5}\\

\ganttgroup{Multi-subtask}{1}{9}\\
\ganttbar{Subtask representation \& data}{1}{3}\\
\ganttbar{Compositional decoder}{3}{6}\\
\ganttbar{Evaluation \& ablations}{6}{9}\\
\ganttmilestone{CVPR 2027 submission}{3}\\
\ganttmilestone{ICCV 2027 submission}{7}\\

\ganttgroup{Thesis}{6}{13}\\
\ganttbar{Chapters 3--5 (from papers)}{6}{8}\\
\ganttbar{Chapters 6--7}{8}{11}\\
\ganttbar{Intro, lit review, conclusion}{10}{12}\\
\ganttbar{Revision}{12}{13}\\
\ganttmilestone{Thesis submission}{13}
\end{ganttchart}
\caption{Plan for the remaining work, September 2026 to September 2027.
Months are shown along the top.}
\label{fig:gantt}
\end{figure}

Chapters~3--5 derive from published or submitted papers and require
reorganisation rather than new writing, so drafting can begin as soon as the
TopX submission is settled. On this plan the thesis is submitted in September
2027. Submission in March--June 2027 is achievable if the multi-subtask work is
taken to a single conference submission rather than developed further; the
September date is preferred because it allows that chapter to be written from an
accepted paper.

\section{Publication plan}

\begin{table}[h]
\centering\small
\begin{tabularx}{\textwidth}{@{}l X@{}}
\toprule
\textbf{Date} & \textbf{Planned work}\\
\midrule
September 2026 & TopX --- ICLR 2027 submission (in progress)\\
November 2026  & Multi-subtask --- CVPR 2027 submission (optimistic)\\
January 2027   & KV-Control decision; camera-ready or resubmission\\
March 2027     & Multi-subtask --- ICCV 2027 submission (realistic)\\
March 2027     & TopX OOD extension --- ICCV 2027 or journal submission\\
May 2027       & Multi-subtask --- NeurIPS 2027 submission (fallback)\\
June 2027      & All paper work complete; writing only from this point\\
July--Sept 2027& Thesis writing and revision\\
September 2027 & Thesis submission\\
\bottomrule
\end{tabularx}
\caption{Planned schedule for the remaining work}
\label{tab:schedule}
\end{table}

\begin{itemize}
\item \textbf{MCRE:} published, ECCV 2024 Workshop on Foundation Models for
  3D Humans.
\item \textbf{Hi-RQCT:} published, CVMP 2025.
\item \textbf{MOGO:} published, AAAI 2026.
\item \textbf{MotionDuet:} published, CVPR 2026.
\item \textbf{KV-Control:} under review at AAAI; fallback submission to CVPR
  2027.
\item \textbf{TopX:} in submission to ICLR 2027; the out-of-distribution
  specialisation work follows as an ICCV 2027 submission or a journal
  extension.
\item \textbf{Multi-subtask motion generation:} CVPR 2027 optimistic, ICCV 2027
  realistic, NeurIPS 2027 fallback.
\end{itemize}

% =================================================================
\bibliographystyle{plain}
\bibliography{references,references_30m}

\end{document}
